\documentclass[10pt]{article}
\usepackage[margin=0.8in,includefoot]{geometry}
\usepackage{longtable,booktabs,array,fancyhdr,hyperref,titlesec}
\newcolumntype{L}[1]{>{\raggedright\arraybackslash}p{#1}}
\pagestyle{fancy}\fancyhf{}
\fancyhead[L]{\sffamily SeaBird Volume 5 -- Binary Interfaces}
\fancyhead[R]{\sffamily Version 3.0 RC1}
\fancyfoot[C]{\thepage}\setlength{\headheight}{14pt}
\titleformat{\section}{\Large\bfseries\sffamily}{\thesection}{0.7em}{}
\titleformat{\subsection}{\large\bfseries\sffamily}{\thesubsection}{0.7em}{}
\hypersetup{colorlinks=true,linkcolor=black}\sloppy\emergencystretch=1em
\begin{document}
\begin{titlepage}\centering\vspace*{2cm}
{\Huge\bfseries SeaBird ISA\\[0.3cm]}{\LARGE Volume 5: Binary Interfaces\\[0.5cm]}
{\Large Version 3.0 Release Candidate 1}\vfill
{\large Generated from architectural-layouts.json\\June 27, 2026}\vfill
\end{titlepage}
\tableofcontents\clearpage
\section{Feature Discovery Bit Assignments}
QUERY leaf 0x0001 returns the basic bitmap in R0. QUERY leaf 0x0003 returns the extended bitmap in R0. Unlisted bits are reserved and read zero.

\subsection{basic\_r0}
\begin{longtable}{@{}L{3.5cm}L{2cm}@{}}
\toprule Feature & Bit \\ \midrule
CLOWNFISH & 0 \\
TETRA & 1 \\
DRAGONET & 2 \\
DROPLET & 3 \\
MMU & 4 \\
ATOMICS & 5 \\
FP & 6 \\
SIMD & 7 \\
PIO & 8 \\
SMP & 9 \\
XSTATE & 10 \\
DEBUG & 11 \\
PMU & 12 \\
MCA & 13 \\
\bottomrule\end{longtable}
\subsection{extended\_r0}
\begin{longtable}{@{}L{3.5cm}L{2cm}@{}}
\toprule Feature & Bit \\ \midrule
AVX & 0 \\
CRYPTO & 1 \\
DSP & 2 \\
TXN & 3 \\
VIRT & 4 \\
SHSTK & 5 \\
IBT & 6 \\
BCD & 7 \\
MLAUNCH & 8 \\
RNG & 9 \\
ATOMIC128 & 10 \\
RINGS12 & 11 \\
\bottomrule\end{longtable}
\section{Control Register Bit Layouts}
\subsection{CR0 (64 bits)}
\begin{longtable}{@{}L{4.5cm}L{2cm}L{2cm}@{}}
\toprule Field & First bit & Width \\ \midrule
PE & 0 & 1 \\
MP & 1 & 1 \\
EM & 2 & 1 \\
TS & 3 & 1 \\
NE & 5 & 1 \\
WP & 16 & 1 \\
AM & 18 & 1 \\
PG & 31 & 1 \\
\bottomrule\end{longtable}
Unlisted bits are reserved, read as zero, and must be written as zero.
\subsection{CR1 (64 bits)}
\begin{longtable}{@{}L{4.5cm}L{2cm}L{2cm}@{}}
\toprule Field & First bit & Width \\ \midrule
FP & 0 & 1 \\
SIMD & 1 & 1 \\
AVX & 2 & 1 \\
CRYPTO & 3 & 1 \\
DSP & 4 & 1 \\
TXN & 5 & 1 \\
PIO & 6 & 1 \\
VIRT & 7 & 1 \\
PMU & 8 & 1 \\
SHSTK & 9 & 1 \\
IBT & 10 & 1 \\
RNG & 11 & 1 \\
\bottomrule\end{longtable}
Unlisted bits are reserved, read as zero, and must be written as zero.
\subsection{CR3 (64 bits)}
\begin{longtable}{@{}L{4.5cm}L{2cm}L{2cm}@{}}
\toprule Field & First bit & Width \\ \midrule
ASID & 0 & 12 \\
ROOT\_PFN & 12 & 40 \\
NO\_FLUSH & 63 & 1 \\
\bottomrule\end{longtable}
Unlisted bits are reserved, read as zero, and must be written as zero.
\subsection{CR4 (64 bits)}
\begin{longtable}{@{}L{4.5cm}L{2cm}L{2cm}@{}}
\toprule Field & First bit & Width \\ \midrule
GLOBAL\_PAGES & 0 & 1 \\
ASID\_ENABLE & 1 & 1 \\
PKEY\_ENABLE & 2 & 1 \\
USER\_ALIGN & 3 & 1 \\
SHSTK\_ENABLE & 4 & 1 \\
IBT\_ENABLE & 5 & 1 \\
XSTATE\_ENABLE & 6 & 1 \\
STAGE2\_ENABLE & 7 & 1 \\
\bottomrule\end{longtable}
Unlisted bits are reserved, read as zero, and must be written as zero.
\subsection{MODE (64 bits)}
\begin{longtable}{@{}L{4.5cm}L{2cm}L{2cm}@{}}
\toprule Field & First bit & Width \\ \midrule
EXECUTION\_MODE & 0 & 2 \\
VA\_LEVELS\_MINUS\_4 & 2 & 1 \\
\bottomrule\end{longtable}
Unlisted bits are reserved, read as zero, and must be written as zero.
\section{Architectural Binary Structures}
All offsets and sizes in this section are bytes. All multi-byte fields are little-endian.
\subsection{IDT\_ENTRY}
Size: 32 bytes. Alignment: 32 bytes.
\begin{longtable}{@{}L{4cm}L{2cm}L{2cm}@{}}
\toprule Field & Offset & Size \\ \midrule
handler\_low & 0 & 8 \\
handler\_high & 8 & 8 \\
target\_cpl & 16 & 1 \\
gate\_type & 17 & 1 \\
stack\_selector & 18 & 1 \\
caller\_cpl & 19 & 1 \\
present & 20 & 1 \\
reserved & 21 & 11 \\
\bottomrule\end{longtable}
\subsection{CORE\_CONTEXT}
Size: 1024 bytes. Alignment: 64 bytes.
\begin{longtable}{@{}L{4cm}L{2cm}L{2cm}@{}}
\toprule Field & Offset & Size \\ \midrule
magic & 0 & 4 \\
revision & 4 & 2 \\
mode & 6 & 1 \\
cpl & 7 & 1 \\
image\_size & 8 & 4 \\
flags & 16 & 16 \\
ip & 32 & 16 \\
gpr\_r0\_r31 & 64 & 512 \\
fpcr & 576 & 8 \\
fpsr & 584 & 8 \\
k0\_k7 & 592 & 64 \\
system\_state & 656 & 112 \\
reserved & 768 & 256 \\
\bottomrule\end{longtable}
\subsection{XSTATE\_HEADER}
Size: 64 bytes. Alignment: 64 bytes.
\begin{longtable}{@{}L{4cm}L{2cm}L{2cm}@{}}
\toprule Field & Offset & Size \\ \midrule
magic & 0 & 4 \\
revision & 4 & 2 \\
header\_size & 6 & 2 \\
total\_size & 8 & 4 \\
present\_mask & 16 & 16 \\
enabled\_mask & 32 & 16 \\
checksum & 48 & 8 \\
reserved & 56 & 8 \\
\bottomrule\end{longtable}
\subsection{VMCB}
Size: 4096 bytes. Alignment: 4096 bytes.
\begin{longtable}{@{}L{4cm}L{2cm}L{2cm}@{}}
\toprule Field & Offset & Size \\ \midrule
magic & 0 & 4 \\
revision & 4 & 2 \\
header\_size & 6 & 2 \\
total\_size & 8 & 4 \\
flags & 16 & 8 \\
stage2\_root & 24 & 8 \\
guest\_asid & 32 & 4 \\
intercept\_bitmap & 64 & 64 \\
guest\_context & 128 & 1024 \\
host\_context & 1152 & 1024 \\
exit\_reason & 2176 & 4 \\
exit\_qualification & 2184 & 8 \\
exit\_gpa & 2192 & 16 \\
injection & 2208 & 32 \\
reserved & 2240 & 1856 \\
\bottomrule\end{longtable}
\section{Extended State Components}
QUERY leaf 0x0005 uses the component number as subleaf. R0 returns component size, R1 alignment, R2 offset in the standard image, and R3 feature requirements. Offsets are ascending, non-overlapping, and aligned.
\begin{longtable}{@{}L{1.5cm}L{3cm}L{2cm}L{5cm}@{}}
\toprule ID & Name & Alignment & Contents \\ \midrule
0 & FP\_MASK & 64 & FPCR, FPSR, and K0-K7 \\
1 & VECTOR & 64 & V0-V31 at enabled vector width \\
2 & DEBUG & 64 & DBGCTL and DB0-DB7 \\
3 & PMU & 64 & enabled performance-monitor state \\
4 & TXN & 64 & transaction status; no live transaction may be saved \\
5 & VIRT & 64 & virtualization host state \\
\bottomrule\end{longtable}
\section{ELF Object ABI}
SeaBird objects are little-endian ELF. Until an official registry assignment exists, experimental tools use e\_machine value 0x5342. Production tools must make the value configurable and must reject incompatible ABI notes.
\subsection{Relocations}
ELF64 SeaBird objects use RELA records. Relocation expressions follow S+A for absolute fields and S+A-P for PC-relative fields, where P is the address of the relocated field. A canonical five-byte B-type rel32 instruction places its field one byte after the opcode and therefore uses addend -4 so the architectural base is next\_IP.
\begin{longtable}{@{}L{4.5cm}L{2cm}@{}}
\toprule Relocation & Number \\ \midrule
R\_SB\_NONE & 0 \\
R\_SB\_ABS16 & 1 \\
R\_SB\_ABS32 & 2 \\
R\_SB\_ABS64 & 3 \\
R\_SB\_ABS128 & 4 \\
R\_SB\_PCREL8 & 5 \\
R\_SB\_PCREL16 & 6 \\
R\_SB\_PCREL32 & 7 \\
R\_SB\_GOT32 & 8 \\
R\_SB\_GOT64 & 9 \\
R\_SB\_PLT32 & 10 \\
R\_SB\_TLS\_GD & 11 \\
R\_SB\_TLS\_IE & 12 \\
R\_SB\_TLS\_LE & 13 \\
R\_SB\_RELATIVE & 14 \\
R\_SB\_JUMP\_SLOT & 15 \\
R\_SB\_GLOB\_DAT & 16 \\
\bottomrule\end{longtable}
\subsection{TLS Models}

TPBASE addresses the current thread control block. Local-exec computes TPBASE plus a link-time constant. Initial-exec loads a TP-relative offset through the GOT. General-dynamic calls the platform TLS resolver with the module and offset pair and receives the address in R0. The resolver follows the standard calling convention and may clobber caller-saved registers.

PLT entries preserve argument registers, use R28--R31 as resolver scratch, and tail-branch to the resolved target. The GOT is pointer-width aligned. Dynamic RELATIVE relocations compute base plus addend without symbol lookup.

\section{DWARF and Unwinding}
The canonical frame address uses register 7 (R7). The return-address column is 32 (IP). The CIE augmentation string is \texttt{zSB}.
\begin{longtable}{@{}L{4cm}L{4cm}@{}}
\toprule Architectural register & DWARF number \\ \midrule
R0-R31 & 0--31 \\
IP & 32 \\
FLAGS & 33 \\
SPTR & 34 \\
TPBASE & 35 \\
KPBASE & 36 \\
FPCR & 37 \\
FPSR & 38 \\
V0-V31 & 64--95 \\
K0-K7 & 96--103 \\
\bottomrule\end{longtable}
\section{Ratification Requirements}

A conforming implementation must pass opcode uniqueness, reserved-encoding, decoder round-trip, arithmetic edge-case, precise-exception, paging-permission, memory-order, context-image, and ABI relocation tests. A reserved instruction must raise INVALID\_OP before privilege, address, or arithmetic checks. Architecture changes require a database revision and regenerated manuals; hand-edited generated files are non-normative.

\section{RC1 Ratification Results}
The RC1 corpus contains 310 normative instructions and 310 ID-bound golden vectors. Exact ID coverage is required. The independent edge suite contains 34 arithmetic, FP, vector, restart, DSP, and cryptographic cases. The Python oracle, C++ reference core, and translator agree on their overlapping byte and execution domains.
\textbf{Result: PASS -- all architecture-owned RC1 ratification gates are satisfied.}
\end{document}
